% ===========================================================================
% Motivating Example -- paper-ready LaTeX
% Pairs with Figure 1 (motivation_figure.tex).
% Replace \ref{sec:corpus}, \ref{sec:taxonomy}, \ref{fig:motivation}.
% No extra packages: \texttt, \textbf, \emph, $..$ are all kernel.
% ===========================================================================
\section{Motivating Example}
\label{sec:motivation}
%
% Local only: lets TeX stretch a little rather than overflow when a long
% unbreakable \texttt{} token (e.g. highlightBubble) lands at the end of a line.
% Verified clean from 8cm to 16cm text width.
\begingroup\emergencystretch=3em

An Agent Skill is a directory an agent loads as instructions: a \texttt{SKILL.md}
contract, together with the scripts and assets that contract refers to. The
agent cannot execute the contract to see whether it holds. It reads the text and
acts on it. We take one Skill from the SkillBench corpus we audit in
Section~\ref{sec:corpus} and read its contract the way an agent would.

\texttt{d3-visualization} is a chart-building Skill. Line~3 of its
\texttt{SKILL.md} describes it as producing visualizations ``with D3.js (v6)''.
Line~75 tells the reader to pin D3 to ``a specific version (default:
\textbf{d3@7.9.0})'', and the next line recommends vendoring that build. The
document names two major versions of its own central dependency, seventy lines
apart.

The same directory ships an example that links a table to the chart. Its header
promises ``Two-way highlighting (click table row $\rightarrow$ highlight chart
element)''. The click handler calls \texttt{highlightBubble}, which
no file in the Skill defines; the chart example defines
\texttt{selectBubble} instead. The call is guarded by a
\texttt{typeof} check, so it raises nothing. It does nothing either.

A third inconsistency runs between the document and the example it ships.
\texttt{SKILL.md:70} requires the layout to ``fix initial positions
deterministically'' and to ``run exactly N ticks and stop''; that same example
builds a default \texttt{d3.forceSimulation}, with neither fixed initial
positions nor a tick limit.

All three are settled by reading the files, and all three are instances of the
inconsistency types we define in Section~\ref{sec:taxonomy}. They are also not
of one kind: a contradiction inside a single document, a mismatch between two
files that must interoperate, and a rule the shipped example violates. Nothing in
the artifact tells a reader which statement to trust
(Figure~\ref{fig:motivation}).

\paragraph{Observation 1.} \emph{A Skill contract can contradict itself, not only
its code.} Two of the defects above are internal to the contract: its own
description and its own pinning rule disagree on the dependency version, and its
own example calls a helper the contract never defines. Neither is settled by
consulting an external fact, and neither involves an implementation file. If such
self-contradictions are common, then treating a Skill as a single trusted
specification is unsafe in a way that documentation-quality checks do not catch.
RQ1 asks how often real Skills contain contract inconsistencies, of which kinds,
and how they are distributed.

\paragraph{Observation 2.} \emph{Being detectable and being consequential are
independent.} All three defects are visible without running anything, yet they
are not interchangeable. The version conflict can be settled by inspecting the
vendored bundle. The undefined call fails silently and costs nothing until
something downstream depends on the highlight. The determinism rule is violated
by the very artifact the contract tells the reader to use. The text gives no
signal about which of the three will change what an agent builds. RQ2 asks
whether inconsistencies can be found automatically, and whether the ones an agent
will act on can be separated from the ones it will read past.

\paragraph{Observation 3.} \emph{Every path through a defective contract costs
something.} An agent that notices the version conflict has to decide which of the
two statements to believe, spending effort the contract should not have created.
An agent that does not notice inherits the conflict into its output, and the
defect surfaces, if at all, far from its cause. These costs are not comparable in
kind, and neither shows up in a pass/fail result. RQ3 asks what an inconsistency
costs an agent, and whether repairing it recovers the cost.
\endgroup
